\clearpage
\VerdierPhysicalPageAnchor{0150}%
\begingroup
\small
\setlength{\parskip}{0pt}%
\setlength{\emergencystretch}{1.5em}%

\VerdierSemanticAnchor{ch2-2-3-5-w-continuation}
\noindent 由于 $B$ 是 $\mathcal B$ 的对象，且 $s$ 属于
$S_{\mathcal D}(\mathcal A\cap\mathcal B)$，对象 $W$ 同构于
$\mathcal B$ 的某个对象 (1.2.4)。
\VerdierSemanticAnchor{ch2-2-3-5-i-implies-ii-conclusion}
因此由第一个断言，态射 $f$ 通过 $\mathcal A\cap\mathcal B$ 的某个对象
分解。于是 $Q_3(f)$ 为零，从而 $m=0$。
\VerdierSemanticAnchor{ch2-2-3-5-ii-implies-i}
现在证明 ii) $\Rightarrow$ i)。由 a)，只需证明：$\mathcal D$ 中的任意态射
$g\colon B\longrightarrow X$，其中 $B$ 是 $\mathcal B$ 的对象且 $X$ 是
$\mathcal A$ 的对象，都通过 $\mathcal A\cap\mathcal B$ 的某个对象分解。
\VerdierSemanticAnchor{ch2-2-3-5-q3-g-zero}
由假设，
$Q_3(g)\in\operatorname{Hom}_{\mathcal D/(\mathcal A\cap\mathcal B)}(B,X)$
为零。
\VerdierSemanticAnchor{ch2-2-3-5-t-morphism}
因此由 (2.2.4)，存在态射 $t\colon X\longrightarrow X'$，其中
$t\in S_{\mathcal D}(\mathcal A\cap\mathcal B)$，且 $tg=0$。
\VerdierSemanticAnchor{ch2-2-3-5-distinguished-triangle}
由 $S_{\mathcal D}(\mathcal A\cap\mathcal B)$ 的定义，$\mathcal D$ 中存在
形如 $(B',X,X',i,t,w)$ 的特异三角，其中 $B'$ 是
$\mathcal A\cap\mathcal B$ 的对象。
\VerdierSemanticAnchor{ch2-2-3-5-proof-complete}
由 (1.2.1)，等式 $tg=0$ 表明态射 $g$ 通过 $B'$ 分解，正合所需。命题
(2.3.5) 得证。
\VerdierSemanticAnchor{ch2-2-3-5-left-localizing-definition}
\medskip
若 $\mathcal D^\circ$ 的子范畴 $\mathcal A^\circ$ 是右
$\mathcal B^\circ$-局部化的，则称 $\mathcal D$ 的满三角子范畴
$\mathcal A$ 是\emph{左 $\mathcal B$-局部化的}。
\endgroup
